%% [AI DRAFT] Borrador generado para discusión (Modo C, etapa 3).
%% Citas verificadas contra el PDF. Sin marcadores [VERIFY] pendientes.

\section{EXPERIMENTOS}

Esta sección presenta un resultado diagnóstico sobre el conjunto de datos del que se
dispone y describe el diseño con el que se someterá a prueba la arquitectura completa. El
primero es un experimento ejecutado; el segundo, un procedimiento previsto. Los umbrales de
éxito de cada métrica no se fijan aquí: dependen de la cobertura de características que arroje
la primera fase, y establecerlos antes de conocerla sería arbitrario.

\subsection{Separabilidad del conjunto disponible}

Antes de evaluar la arquitectura conviene establecer qué puede concluirse del conjunto
disponible. Se trata de un conjunto \textbf{sintético} de 10\,000 registros con tres clases
equilibradas, construido ante la ausencia de un conjunto etiquetado público de fichas de
vendedores peruanos: sus variables y rangos se derivaron de los factores de riesgo
identificados en la Sección~II, y las etiquetas se asignaron mediante reglas deterministas
sobre esas mismas variables. Declarar esta procedencia es indispensable para interpretar lo
que sigue. Sus siete entradas
corresponden a siete de las ocho de la Tabla~\ref{tab:senales}: la veracidad de dominio, que
este trabajo introduce, es la única que no forma parte de ese vector. Se entrenaron cinco algoritmos de familias
distintas sobre la misma partición 80/20 con semilla fija, y se validaron además mediante
validación cruzada estratificada de cinco pliegues sobre el conjunto completo. La
Tabla~\ref{tab:bench} recoge los resultados.

Los experimentos se ejecutaron en un entorno estandarizado para garantizar su
reproducibilidad. A nivel de software se utilizó Ubuntu 24.04.3 LTS (kernel 6.8.0-137-generic)
con Python 3.11 y las bibliotecas scikit-learn, NumPy y pandas. El hardware consistió en un servidor
Supermicro equipado con un procesador Intel Xeon E3-1240 v5 de cuatro núcleos y ocho hilos a
3.90\,GHz, con los procesos ejecutados íntegramente en CPU y sin aceleración gráfica dedicada.
Las máquinas de soporte vectorial y el clasificador de vecinos más cercanos se encadenaron a
un normalizador de media cero y varianza unitaria, por ser los únicos algoritmos del conjunto
sensibles a la escala de las características. Los demás hiperparámetros se mantuvieron en sus
valores por defecto, con semilla fija en todos los casos: el objetivo no era optimizar el
desempeño sino comprobar si existía diferencia alguna entre familias de algoritmos.

\begin{table}[t]
\caption{Desempeño de cinco algoritmos sobre el conjunto disponible}
\label{tab:bench}
\centering
\footnotesize
\renewcommand{\arraystretch}{1.12}
\begin{tabular}{@{}p{2.55cm}ccc@{}}
\toprule
\textbf{Algoritmo} & \textbf{Exactitud} & \textbf{F1 macro} & \textbf{CV (5)} \\
\midrule
Bosque aleatorio     & 1.0000 & 1.0000 & 1.0000 $\pm$ 0.0000 \\
XGBoost              & 1.0000 & 1.0000 & 1.0000 $\pm$ 0.0000 \\
Árbol de decisión    & 1.0000 & 1.0000 & 1.0000 $\pm$ 0.0000 \\
SVM (RBF)            & 1.0000 & 1.0000 & 1.0000 $\pm$ 0.0000 \\
$k$-NN ($k=5$)       & 1.0000 & 1.0000 & 1.0000 $\pm$ 0.0000 \\
\bottomrule
\end{tabular}
\end{table}

Los cinco algoritmos clasifican correctamente los 2\,000 registros del conjunto de prueba, con
desviación nula entre pliegues. El resultado no debe leerse como evidencia de la calidad de
ninguno de ellos. Que un clasificador de vecinos más cercanos iguale a un ensamble de
gradiente potenciado, y que ambos igualen a un árbol individual, indica que la frontera de
decisión es recuperable por cualquier método: el problema es separable por construcción.

Un segundo análisis lo confirma. Un árbol de decisión de profundidad cinco reproduce la regla
de etiquetado en su totalidad, y la importancia de las características se concentra en dos (la
puntuación del producto, con 0.51, y el volumen de ventas, con 0.46), mientras que tres de las
siete no contribuyen en absoluto. Entre estas últimas figura la proporción estimada de reseñas
fraudulentas, que es precisamente una de las características sobre las que descansa el
argumento de la Sección~IV-E.

De ello se siguen dos consecuencias. La primera es que este conjunto \textbf{no permite
discriminar entre algoritmos}: la elección del bosque aleatorio descrita en la Sección~IV-E se sostiene
en sus propiedades de interpretabilidad y costo, no en una comparación empírica que estos
datos no pueden arbitrar. La segunda, más relevante, es que las métricas obtenidas sobre él no
estiman la calidad del nivel de riesgo asignado, sino la fidelidad con que un modelo recupera
la regla que generó sus propias etiquetas. La construcción de un conjunto con
etiquetas de origen independiente del vector de entrada es, por tanto, condición previa a
cualquier afirmación sobre eficacia.


\subsection{Diseño experimental de la validación}

La validación se ejecutará en el Perú, sobre las plataformas y los hábitos descritos en la
Sección~III-A. Esa elección delimita el alcance de las conclusiones empíricas, no el de la
arquitectura: las características de la Tabla~\ref{tab:senales} se derivan de elementos
presentes en cualquier ficha de producto y no de particularidades del mercado peruano. La
propuesta es trasladable a otros contextos, con la salvedad de que la cobertura efectiva de
cada característica debe volver a medirse en cada mercado.

La validación se organiza en dos fases con objetivos separables.

\emph{Fase 1: observabilidad de las características.} Antes de evaluar la utilidad del veredicto es
necesario establecer qué fracción de las características de la Tabla~\ref{tab:senales} resulta
efectivamente legible en fichas reales. Sobre una muestra de fichas de producto de las
plataformas de mayor uso en el Perú se registrará, para cada característica, si pudo resolverse o si
quedó indeterminada, distinguiendo entre plataformas grandes y vendedores independientes.
El desenlace relevante no es un porcentaje agregado sino su distribución: una arquitectura
que solo resolviera características en las plataformas que ya ofrecen garantías carecería de
utilidad precisamente en el segmento donde el riesgo se concentra.

\emph{Fase 2: eficacia con usuarios.} Un diseño con dos grupos, con y sin la extensión activa,
enfrentará a los participantes a un conjunto de fichas preparadas que incluya casos legítimos y
casos con indicadores de fraude. La población objetivo son compradores digitales peruanos de
entre 18 y 60 años, con una muestra prevista de cincuenta participantes que otorgarán su
consentimiento de forma voluntaria antes de comenzar. Las sesiones se conducirán de manera
presencial o en línea, según la disponibilidad de cada participante. Los registros de consulta
descritos en la Sección~IV-F se recogerán disociados del participante que los originó, se
conservarán únicamente durante el periodo de análisis y se eliminarán al cierre del estudio;
ningún dato que permita reconstruir el historial de compra de una persona identificada forma
parte de lo que el protocolo almacena.

Del comportamiento observado durante la sesión se medirá la tasa de identificación correcta,
la tasa de falsos positivos sobre comercios legítimos y el tiempo hasta la decisión. La
percepción del participante se recoge mediante un cuestionario estructurado de quince ítems
con respuesta en escala de Likert de cinco puntos, organizado en las cinco dimensiones de la
Tabla~\ref{tab:instrumento}; el bloque de usabilidad sigue la Escala de Usabilidad del Sistema
y se informa sobre su escala de 0 a 100 \cite{ref23}. Declarar la composición del instrumento
permite situar cada conclusión en la dimensión que efectivamente la sostiene: la eficacia se
juzga por las tres primeras medidas de comportamiento, y la aceptación, por las dos últimas
dimensiones del cuestionario.

\begin{table}[t]
\caption{Dimensiones del instrumento de validación con usuarios}
\label{tab:instrumento}
\centering
\footnotesize
\setlength{\tabcolsep}{3.5pt}
\renewcommand{\arraystretch}{1.15}
\begin{tabular}{@{}p{2.35cm}p{5.3cm}@{}}
\toprule
\textbf{Dimensión} & \textbf{Qué establece} \\
\midrule
Usabilidad & Si el veredicto y su explicación se interpretan sin instrucción previa (bloque SUS) \\
Confianza percibida & Si el participante considera fundada la advertencia que recibe \\
Riesgo percibido & Si su estimación del riesgo de la ficha varía al disponer de la herramienta \\
Intención de uso y recomendación & Si declararía seguir usándola fuera del entorno de prueba \\
Influencia en la decisión & Si la advertencia modificó lo que habría hecho sin ella \\
\bottomrule
\end{tabular}
\end{table}

La tasa de falsos positivos es la métrica que gobierna el diseño. Un instrumento que
advierte con excesiva frecuencia pierde credibilidad y deja de ser consultado, y su costo recae además sobre
comercios honestos que no publicitan garantías. Por ello la regla de abstención de la
Sección~IV distingue el veredicto de riesgo del de evidencia insuficiente: son errores
distintos y deben medirse por separado.

\subsection{Amenazas a la validez}

Cuatro limitaciones condicionan lo que este diseño podrá concluir, y conviene explicitarlas
antes de su ejecución.

La primera afecta a la capa cliente. La extracción depende de selectores sobre el DOM, y se
ha advertido que las técnicas apoyadas en selectores predefinidos son sensibles a cambios de
maquetación \cite{ref24}. Una plataforma
que rediseñe su ficha reducirá la cobertura de características sin previo aviso, por lo que la
Fase~1 debe repetirse periódicamente y no tratarse como una medición única.

La segunda afecta al clasificador de riesgo. Sus etiquetas de entrenamiento se derivan
mediante reglas deterministas aplicadas a las mismas características que el modelo recibe como
entrada, de modo que el desempeño medido sobre ese conjunto refleja la recuperación de esa
regla y no el comportamiento frente a vendedores fraudulentos reales. Su evaluación exige, por
tanto, un conjunto de datos con etiquetas de origen independiente del vector de
características.

La tercera es de validez ecológica. Un experimento con fichas preparadas mide la capacidad de
detección en condiciones de atención deliberada, mientras que la evidencia disponible sugiere
que la victimización ocurre justamente cuando esa atención falla \cite{ref27}. El diseño
puede establecer si la herramienta ayuda a quien la mira, no si protege a quien compra
distraído.

La cuarta es la más general: la arquitectura supone que el atacante no adapta su
comportamiento a las características observadas, suposición insostenible a mediano plazo
porque los defraudadores ajustan sus estrategias frente a los mecanismos de detección
\cite{ref28}. El registro de consultas de la Sección~IV-F existe, entre otras razones, para
hacer observable esa deriva.
